\documentclass[a4paper,12pt]{article}
\usepackage[utf8]{inputenc}
\usepackage[T1]{fontenc}
\usepackage[ngerman]{babel}
\usepackage{geometry}
\geometry{margin=2.5cm}
\usepackage{amsmath}
\usepackage{parskip}
\setlength{\parskip}{6pt}

\begin{document}

\section*{Toward a Better Understanding of Fourier Neural Operators from a Spectral Perspective}

\textbf{Autoren:} Shaoxiang Qin, Fuyuan Lyu, Wenhui Peng, Dingyang Geng, Ju Wang, Xing Tang,
Sylvie Leroyer, Naiping Gao, Xue Liu, Liangzhu (Leon) Wang \\
\textbf{Jahr:} 2024 \\
\textbf{Konferenz/Preprint:} arXiv:2404.07200v2

% -------------------------------------------------------
\subsection*{Zusammenfassung}

Fourier Neural Operators (FNOs) haben sich als besonders effektive Surrogatmodelle f\"ur das
L\"osen partieller Differentialgleichungen (PDEs) etabliert. Trotz ihrer vielfach demonstrierten
Leistungsf\"ahigkeit besteht ein bislang unzureichend verstandenes Problem: FNOs profitieren
kaum von gr\"o\ss{}eren Fourier-Kernen, die ein breiteres Frequenzspektrum abdecken. Die
Standardl\"osung in der Literatur war es, die Gr\"o\ss{}e des Fourier-Kerns (den sogenannten
Truncation-Mode $k$) m\"oglichst klein zu halten -- was jedoch die F\"ahigkeit des Modells
einschr\"ankt, komplexe PDE-Daten zu repr\"asentieren.

Diese Arbeit liefert eine systematische spektralanalytische Untersuchung des FNO und identifiziert
die Ursache dieses Problems als \textit{Fourier Parameterization Bias}: Konvolutionskerne, die
im Fourierraum parametrisiert werden, zeigen eine deutlich st\"arkere Tendenz, dominante
Frequenzen der Zieldaten zu lernen, w\"ahrend nicht-dominante Frequenzen vernachl\"assigt
werden. Konkret bedeutet dies: Da die Energie von PDE-Daten im Fourierraum typischerweise
stark auf wenige dominante Frequenzmoden konzentriert ist (im Navier-Stokes-Beispiel der
Autoren enthalten $1{,}2\,\%$ der Pixel mit den gr\"o\ss{}ten Merkmalen $99\,\%$ der Gesamtenergie),
optimiert die Verlustfunktion w\"ahrend des Trainings haupts\"achlich diese wenigen Moden.
Zus\"atzliche Frequenzparameter, die durch einen vergr\"o\ss{}erten Kern hinzukommen, werden
kaum genutzt und erzeugen stattdessen Rauschen.

Zur Illustration werden sogenannte NMSE-Spektren eingesetzt: Die Vorhersageresiduen werden
in den Fourierraum transformiert, und die Energieverteilung \"uber Frequenzmoden wird
analysiert. Verglichen mit Conv-basierten Modellen wie U-Net, das Kerne im Ortsraum
parametrisiert, zeigt FNO unterhalb seiner Abschneidefrequenz eine stark \"uberproportionale
Reduktion des Fehlers bei den dominanten Frequenzen, w\"ahrend es bei nicht-dominanten
Frequenzen schlechter abschneidet.

Als L\"osungsansatz schlagen die Autoren \textbf{SpecB-FNO} (Spectral Boosted FNO) vor.
Die zentrale Beobachtung ist, dass die Residuen eines trainierten FNO eine gleichm\"a\ss{}igere
Energieverteilung im Fourierraum aufweisen als die Originaldaten, da der FNO die dominanten
Frequenzen bereits gut erfasst hat. SpecB-FNO nutzt diesen Umstand durch iteratives
Residuallernen: Nach dem Training eines Basis-FNO $\mathcal{G}_0$ werden $T$ zus\"atzliche
FNO-Module sequenziell trainiert, wobei jedes Modul $\mathcal{G}_i$ das Vorhersageresiduum
des vorherigen Moduls lernt:
\[
    r_i = y - \sum_{j=0}^{i-1} \hat{r}_j, \qquad \hat{r}_i = \mathcal{G}_i\!\left([x,\, \hat{r}_{i-1}]\right).
\]
Die finale Vorhersage ergibt sich als Summe aller Modul-Ausgaben:
$\hat{y} = \sum_{i=0}^{T} \hat{r}_i$.

Die experimentelle Evaluation auf f\"unf PDE-Datens\"atzen (Navier-Stokes, Darcy Flow,
Flachwassergleichung, Diffusions-Reaktions-Gleichung) zeigt, dass SpecB-FNO den
Vorhersagefehler gegen\"uber dem besten Baseline-Modell im Durchschnitt um $50\,\%$ reduziert,
in einzelnen F\"allen sogar bis zu $93\,\%$. Dar\"uber hinaus erweist sich SpecB-FNO als
speichereffizient, da stets nur ein Teilmodell trainiert wird, was den maximalen GPU-Speicherbedarf
gegen\"uber vergleichbaren vergr\"o\ss{}erten FNO-Varianten deutlich senkt.

% -------------------------------------------------------
\subsection*{Bezug zur Masterarbeit}

Die Erkenntnisse dieser Arbeit sind auf mehreren Ebenen direkt relevant f\"ur die vorliegende
Masterarbeit, in der FNO als Surrogatmodell zur Vorhersage von Str\"omungsfeldern um
sedimentierende Kristalle im Stokes-Regime untersucht wird.

\textbf{Frequenzinhalt von Stokes-Str\"omungsfeldern.}
Im Stokes-Regime sind Str\"omungsfelder um Kristalle in der Regel glatt und werden von
niederfrequenten Komponenten dominiert -- insbesondere durch die weitreichenden
Druckgradienten und Wirbelstrukturen. An den Kristallr\"andern k\"onnen jedoch starke
lokale Gradienten auftreten, die h\"ohere Frequenzanteile erzeugen. Der \textit{Fourier
Parameterization Bias} liefert eine plausible Erkl\"arung daf\"ur, warum ein Standard-FNO
diese randnahen Details schlechter erfasst als die grob-skaligen Str\"omungsstrukturen --
ein Verhalten, das in den Evaluationsmetriken der Masterarbeit als r\"aumlich ungleichm\"a\ss{}iger
Fehler sichtbar werden kann.

\textbf{Wahl des Truncation-Mode $k$.}
Ein zentrales Hyperparameter-Design des FNO ist die Anzahl der beibehaltenen Fourier-Moden.
Qin et al. zeigen empirisch, dass eine naive Vergr\"o\ss{}erung von $k$ die Genauigkeit nicht
verbessert und sogar verschlechtern kann. F\"ur die Masterarbeit impliziert dies, dass bei
der Hyperparameteroptimierung des FNO nicht automatisch ein m\"oglichst gro\ss{}es $k$ gew\"ahlt
werden sollte; stattdessen sollte $k$ an die tats\"achlich vorliegende Frequenzstruktur der
Str\"omungsfelder angepasst werden, die mit Methoden wie dem NMSE-Spektrum charakterisiert
werden kann.

\textbf{Residuallernen als Erweiterungsm\"oglichkeit.}
Der SpecB-FNO-Ansatz ist konzeptuell verwandt mit dem in der Masterarbeit untersuchten
Residuallernansatz ($v_{\text{total}} = v_{\text{Stokes}} + \Delta v_{\text{learned}}$),
bei dem das Netzwerk Korrekturen zur analytischen Stokes-L\"osung lernt. W\"ahrend der
Residualansatz der Masterarbeit physikalisch motiviert ist und die analytische L\"osung als
Vorwissen nutzt, adressiert SpecB-FNO das spektrale Residual iterativ ohne Vorannahmen
\"uber die Physik. Eine Kombination beider Ideen -- physikalisch informiertes Residuallernen
mit anschlie\ss{}endem spektralem Boosting -- k\"onnte eine interessante Erweiterung darstellen.

\textbf{Vergleich FNO vs.\ U-Net.}
Die Arbeit zeigt quantitativ, dass FNO und U-Net komplementäre St\"arken besitzen: FNO
\"ubertrifft U-Net bei dominanten (niederfrequenten) Strukturen, w\"ahrend U-Net eine
konsistentere spektrale Performance \"uber alle Frequenzen zeigt. Da die Masterarbeit
beide Architekturen vergleicht, bieten die NMSE-Spektren eine zus\"atzliche Analysedimension
jenseits globaler Fehlermetriken: Ein Vergleich der spektralen Fehlerverteilungen kann
erkl\"aren, in welchen r\"aumlichen Skalen U-Net oder FNO jeweils \"uberlegen ist und ob die
Unterschiede auf den Fourier Parameterization Bias zur\"uckzuf\"uhren sind.

\textbf{Generalisierung auf mehrere Kristalle.}
Die Kernfragestellung der Masterarbeit -- ob ein auf $N$ Kristallen trainiertes Modell auf
$N{+}m$ Kristalle generalisiert -- hat eine direkte spektrale Dimension: Mehr Kristalle
erzeugen komplexere Interferenzmuster in den Str\"omungsfeldern, was den Anteil
nicht-dominanter Frequenzen vergr\"o\ss{}ert. Der Fourier Parameterization Bias k\"onnte
demnach bei Multi-Kristall-Konfigurationen st\"arker zum Tragen kommen als bei einfachen
Einzel-Kristall-Systemen, was m\"oglicherweise die beobachtete schlechtere Generalisierung
des FNO auf komplexere Konfigurationen erkl\"art.

\end{document}